\section{Scope and Evidence Method}
\label{sec:scope}
\paragraph{Unit of analysis.}
We consider work addressing composition-level layout or graphics generation, design-specific visual text, construction of structured or layered assets, editing or reconstruction of graphic designs, and evaluation of these tasks. Domains include posters, advertisements, infographics, magazine pages, scientific figures, and slides. Generic image generation is included only if a graphic-design-specific constraint, structure, or evaluation is central.

\paragraph{Discovery and selection.}
The working evidence register is seeded from a public research catalogue, prior surveys, original research proceedings, author project pages, and references within candidate papers. A publication and its preprint versions should be deduplicated as one research contribution after verifying their relationship. The catalogue is a living research index, not a final screened sample. We code multi-task papers on multiple analytical axes, even if they have only one primary catalogue category.

\paragraph{Verification levels.}
We distinguish \emph{catalogue-only} entries, \emph{author/publisher abstract or indexed} entries, and works for which at least selected \emph{original methods or results} have been inspected. The last group is not necessarily a full-text audit of all appendices or an independent replication. Empirical statements require identifiable evidence such as a method section, table, figure, or publisher record. Missing evaluation details in an uninspected paper are not treated as evidence that no evaluation exists.

At the November 3, 2026 audit checkpoint, the broader provisional register held 225 entries: 88 with selected original method/result sections checked, 28 supported through abstract or publication-index evidence, and 109 catalogue-only. These figures describe \emph{verification depth}, not a PRISMA screening outcome, a claim that all 225 studies were comprehensively read, or the size of the stable public catalogue.

\paragraph{Coding protocol.}
For each work, we record input and task, output representation, control mechanism, editing operation, measurement target, metric/protocol, and source depth. Boundaries matter: SVG paths are not necessarily editable semantic text; raster inpainting is a real image-editing capability without implying Photoshop-style object layers; and a graphical user interface can be demonstrated without a controlled user study.

\paragraph{Limitations of the review process.}
This is an initial narrative survey with structured coding, \emph{not yet} a reproducible systematic review. Search strings, database query dates, deduplication logs, exclusion rules, independent screening decisions, and reviewer agreement must be finalized before presenting complete coverage. The bibliography is deliberately representative rather than a claim to cite every study in the current registry. In the remainder, we avoid pooling incomparable metrics into a single leaderboard.
